\documentclass[12pt]{article}

\usepackage{amsmath}
\usepackage{graphicx}
\usepackage{natbib}
\usepackage{german}
\DeclareMathOperator\sgn{sgn}

%\graphicspath{{/home/ritz/Teaching/fcn/fig/}}

%----------------- layout ----------------------------------------

\oddsidemargin=0in
\textwidth=6in
\topmargin=-0.7in
\textheight=9.5in
\baselineskip=16pt
\font\mittel=cmr17
\font\roman=cmr12
\font\klein=cmr9
\font\mgross=cmr10
\newfont{\CAP}{cmcsc10 scaled 900}
\pagestyle{empty}
\parindent=0mm

%----------------- macros ----------------------------------------
\newcommand{\conc}[1]{[{\rm C}]_{\text{#1}}}
\newcommand{\units}[1]{\text{#1}}
\newcommand{\K}{\text{K}}
\newcommand{\Na}{\text{Na}}
\newcommand{\Cl}{\text{Cl}}
\newcommand{\Leak}{\text{leak}}
\newcommand{\water}{\text{H}_2\text{O}}
\newcommand{\tderiv}[1]{\frac{\mathrm{d}#1}{\mathrm{d}t}}
%-----------------------------------------------------------------

\renewcommand{\refname}{Literature}

\begin{document}

\parbox{2cm}{
\includegraphics[width=1.8cm]{../bccnlogo.pdf}
}
\parbox{11cm}{
\begin{center}
\large HUMBOLDT-UNIVERSIT"AT \hskip 0.1 cm ZU \hskip 0.1 cm BERLIN
\vskip 0.1 true cm
\mgross BERNSTEIN CENTER FOR\\ COMPUTATIONAL NEUROSCIENCE
\end{center}
}
\parbox{2cm}
{
\hfill
\includegraphics[width=1.8cm]{../hublogo.pdf}
}

\vskip 0.6 true cm

\leftline{\CAP Humboldt-Universit"at zu Berlin 
\hfill Phone: 030/2093-9110}
\leftline{\CAP Philippstr. 13 House 6
             \hfill Fax: 030/2093-6771}
\leftline{\CAP \hfill webpage: http://www.bccn-berlin.de/}


\vskip 0.8 true cm
\centerline{\bf Models of Neural Systems, WS 2010/11}
\centerline{\bf Project 10: Short-term synaptic plasticity} 
\centerline{Project presentation and report submission: February,
15th, 2011}

\vskip 0.8 true cm

\begin{enumerate}
    \item \textbf{Tsodyks-Markram model of short-term synaptic
        depression}        

        \begin{enumerate}
            \item Implement the model of short-term synaptic
                depression by \citet{Tsodyks97}.
            \item Using short spike trains with appropriately-chosen inter-spike
                intervals show the effects of depression and recovery.
            \item Demonstrate how depression and recovery depend on
                parameters of the model: $\tau_{\mathrm{rec}}$,
                $U_{\mathrm{SE}}$ and $\tau_{\mathrm{inact}}$.
        \end{enumerate}

    \item \textbf{Limiting frequency}

        \begin{enumerate}
            \item Feed the dynamic synapse with an input modelled with
                a Poisson spike train with a pre-defined firing rate
                and no refractoriness.
            \item Plot the steady-state amplitude of excitatory post-synaptic
                current  (EPSC) as a function of the input firing
                rate.
            \item Demonstrate that over a certain input frequency
                (limiting frequency) the output firing rate saturates
                and transmits no more information about the input firing
                rate.
        \end{enumerate}

    \item \textbf{Regulating neural code}

        \begin{enumerate}
            \item Simulate the responses of the neuron to inputs with
            constant firing rates and fast transients.
        \item Investigate in which parameter regimes the synapse
            transmits information about either mean firing or  the
            transients.
        \end{enumerate}

    \item \textbf{Gain modulation}

        This part of the project may be solved in one of simulator of
        spiking neural networks (PyNN or Brian).

The goal of this part is to show that synaptic depression may
underline multiplication of neuronal signals \citep{Rothman09}.

\begin{enumerate}
\item Simulate a leaky integrate-and-fire neuron with inhibitory and
    excitatory synapses, where only the excitatory synapses are
    depressing.

\item Plot an output firing rate as a function of the input firing
    rate at different levels of inhibition. Show that inhibition
    modulates the input-output gain.

\item Repeat the above experiment without synaptic depression in the
    excitatory synapses. How
    does the input-output relation change?
\end{enumerate}


\end{enumerate}


\begin{thebibliography}{9}

    \bibitem[Tsodyks and Markram (1997)]{Tsodyks97}
      M. Tsodyks and H. Markram. \emph{The neural code between neocortical
      pyramidal neurons depends on neurotransmitter release
      probability.} Proceedings of the National Academy of Sciences
      94 (2): 719-723, 1997

  \bibitem[Rothman et al. (2009)]{Rothman09}
       J.S. Rothman, L. Cathala, V. Steuber  and R.A. Silver.
       \emph{Synaptic depression enables neuronal gain control.} Nature
       457 (7232): 1015--1018, 2009
      \end{thebibliography}


\vfill
\centerline{\CAP Contact} \CAP

\begin{tabular}{lll}
Urs Bergmann & Phone: 2093-8924 & Email:
urs.bergmann(at)biologie.hu-berlin.de \\
Matthias Guggenmos & & Email: matthias.guggenmos(at)bccn-berlin.de \\
Richard Kempter \hfill & Phone: 2093-8925 \hfill & Email:
r.kempter(at)biologie.hu-berlin.de \\
Bartosz Telenczuk & Phone: 2093-8838 &
Email: b.telenczuk@biologie.hu-berlin.de \\
\end{tabular}

\end{document}

\bibliographystyle{plainnat}
